\documentclass[11pt,a4paper]{article}

\usepackage[LGR,T1]{fontenc}
\usepackage[utf8]{inputenc}
% Transliteraciones con macron+agudo (megalḗtōr, kynṓpidos...)
\DeclareUnicodeCharacter{1E17}{\'{\=e}}
\DeclareUnicodeCharacter{1E53}{\'{\=o}}
\usepackage[greek.ancient,spanish]{babel}
\usepackage[a4paper,margin=2.8cm]{geometry}
\usepackage{microtype}
\usepackage[hidelinks]{hyperref}

% Griego politónico escrito directamente en UTF-8:
% \gr{...} para palabras sueltas; bloques con \foreignlanguage{greek}{...}
\newcommand{\gr}[1]{\foreignlanguage{greek}{#1}}

\hypersetup{
  pdftitle={Leer la Odisea en tiempos iletrados. Sexta parte: Esqueria}
}

\setlength{\parindent}{1.25em}
\setlength{\parskip}{0.35em}
\raggedbottom

\title{Leer la \emph{Odisea} en tiempos iletrados \\ Sexta parte \\ Esqueria}
\author{}
\date{}

\begin{document}

\maketitle

\section*{El palacio de Alcínoo}

Cuando por fin Odiseo se decide a salir del bosque y dirigirse a palacio, Atenea se deja de remilgos y vuelve a la carga ---esta mujer no se está nunca quieta---. Primero se hace la encontradiza, fingiéndose una niña, y luego guía al héroe hasta la residencia de Alcínoo:

\begin{verse}
A la que dijo ella así, por delante fue Palas Atena\\
rauda guiándolo, y él tras la diosa pisando sus huellas.\\
Los navicélebres feacios tampoco se daban ni cuenta\\
de que anduviera por la ciudad entre ellos; ni Atena\\
lo permitía, la diosa zahorí, trenzabella, que en niebla\\
de maravilla envolvía al ser de su benevolencia.\\
Y, boquiabierto, Odiseo los puertos, las naos niveleñas,\\
las tantas plazas de próceres, y las murallas contempla,\\
largas, altivas, con picas ---asombro de ver--- de acroteras.\\
Cuando a la muy alta casa del alto señor ellos llegan,\\
así la diosa, ojiglauca Atenea, le abría conversa:
\end{verse}

Observa, atenta lectora, exacto lector, cómo Homero nos dibuja, en unas pocas líneas, la próspera y marina tierra de los \emph{navicélebres} feacios. Los puertos, las plazas, las naos niveleñas, las largas y altivas murallas. El exiliado, boquiabierto, comprende que ha llegado a un lugar civilizado, quizás el lugar más civilizado que ha visitado nunca. Si Nausícaa puede simbolizar el amor imposible, Esqueria, la tierra de los feacios, representa, creo, el paraíso perdido.

A la puerta de palacio, Atenea, todavía disfrazada, le larga un tremendo discurso a Odiseo informándole de toda la genealogía de Alcínoo y su esposa Arete antes de marcharse. Nuestro héroe entra en palacio y se maravilla de todas las riquezas que encuentra. Homero no escatima detalles:

\begin{verse}
Que era el fulgor de la luna o del sol lo que allí relumbraba\\
por esas salas altivas de Alcínoo buenasentrañas.\\
Y es que paredes de bronce corrían a banda y a banda\\
desde la entrada hasta el fondo, en un friso de azul rematadas;\\
de oro el portón que cerraba por dentro la sólida casa,\\
y las sus jambas de plata, al suelo bronceño clavadas,\\
de plata encima el dintel, y oridorada la aldaba.
\end{verse}

La descripción de la fabulosa morada continúa por un buen rato y uno tiene la sospecha de que la casa de Ulises en Ítaca e incluso el palacio de Menelao en Esparta son chozas cabreras, comparadas con todo el esplendor que al náufrago se le ofrece. El Poeta, que a contable no le gana nadie cuando quiere ponerse a ello, no solo da una descripción minuciosa de ornamentos, oros, platas y bronces, sino que también hace recuento de esclavos, sirvientes, telas, perales, granados, manzanos, higueras, olivos y viñedos que dan exquisito vino. Jauja, en una palabra.

El caso es que el héroe, siempre envuelto en niebla, se las compone para llegar al trono y allí abraza las rodillas de Arete y le dice:

\begin{verse}
«Hija del pardelosdioses Rexénor, Arete, de hinojos\\
a tu señor, tus rodillas, tras tanto pasado, me postro\\
y a los aquí convidados. Vivir os dé el cielo gozosos,\\
y que entregar en legado a los hijos podáis un día todos\\
de vuestra casa el caudal y el honor que os dio el pueblo en decoro.\\
Aparejadme una escolta que a casa me lleve, os lo imploro,\\
pronto, que ha mucho que sufro, y estoy sin los míos y solo».
\end{verse}

Fíjate bien, compasiva lectora, perdonatodo lector. Este hombre que ha saqueado Troya ---y no contento con ello intentó robar también a los cícones--- se ve reducido a la condición de mendigo, se ve obligado a implorar donde antes se habría limitado a arrebatar. En Esqueria descubre ---¿o finge?--- la humildad, aprende a suplicar, ve, quizás por vez primera, los toros de puntiagudos cuernos desde el otro lado de la barrera. Esqueria no solo es la quintaesencia de la hospitalidad. Es el lugar donde Odiseo se vuelve algo más humano.

Y Nolan la elimina de su historia.

Después de abrazar las rodillas de la reina, Ulises se dirige al hogar, donde se sienta «sobre las cenizas, cerca del fuego».
El viejo Equeneo le dice entonces al rey:

\begin{verse}
«Alcínoo, no es para ti ni nobleza ni bien que ese huésped\\
ante el fogón, por sobre la ceniza, en el suelo se siente;\\
y, porque esperan tu voz, todos estos aquí se contienen.\\
Haz levantar a ese hombre y asiento en silla tú ofrécele\\
claveteada de plata, y manda también que nos mezclen\\
más vino aún los heraldos, que ahora al que juega a que truene\\
le arramaremos, a Zeus, que acompaña al que en buen ruego viene;\\
y al extranjero comida dé el ama de lo que almacene».
\end{verse}

El buen rey no se hace de rogar, recoge al mendigo, lo lleva a su mesa, lo alimenta y lo agasaja. Luego despide a todo el mundo, no sin avisar de que al día siguiente planearán el regreso del huésped. Y aquí aparecen las célebres líneas que justifican formalmente la \gr{ξενία}, esa hospitalidad debida a los extraños:

\begin{verse}
Y si resulta que él es un sin muerte que el cielo aquí manda,\\
señal será de que el cielo alguna otra cosa se trama.\\
Que ante nosotros los dioses han hecho visible su estampa\\
siempre que se hizo en su honor ciemboyada gloriosa, y es fama\\
que a nuestro lado han comido sentándose en nuestras bancadas.
\end{verse}

Odiseo, por su parte, responde:

\begin{verse}
«Sea otro, Alcínoo, el cuidado en que des en pensar; no me veo\\
yo como aquellos sin muerte que pueblan la campa del cielo\\
ni en porte ni en condición, y sí un hombre de los morideros.\\
A ese cualquiera que veis soportar de más bregas el peso,\\
a ese hombre sí soy igual en cosa de sufrimiento.\\
Y hasta podría yo hacer de desdichas más grande recuento,\\
cuantas ---y bien que lo sé--- por querer de los dioses yo llevo.\\
Mas permitid que ahora cene, dolido y todo que vengo;\\
cosa más perra no hay que ese impúdico estómago nuestro\\
que nos reclama y apremia a tenerlo a la fuerza en recuerdo\\
por agobiados que estemos y penas haya en el pecho;\\
sí, que mil penas yo llevo en el pecho, y él anda pidiendo\\
siempre que coma y que beba, y que ya no me acuerde de eso,\\
de lo que tanto he pasado, y me llama a dejarlo bien lleno.\\
Pero vosotros, del alba a las claras, haced ---os lo ruego---\\
por prepararle a este pobre infeliz a su patria el regreso,\\
con sus pesares a cuestas; que no se me marche mi tiempo\\
sin ver mi hacienda otra vez, la casa techalta y los siervos».
\end{verse}

Sin desperdicio la reflexión que se marca sobre las exigencias del estómago. Tampoco lo tiene oírle llamarse «pobre infeliz».

Cuando todo el mundo se marcha y se quedan solos, Arete, muy aguda, se fija en las ropas que lleva y deduce, sin decirlo, que se las ha prestado Nausícaa. Dice entonces:

\begin{verse}
«Bien, lo primero, extranjero, te pido yo misma respuesta:\\
¿quién y de qué gentes tú? ¿Quién te ha dado esas ropas que llevas?\\
¿No dices tú que del mar, errabundo infeliz, aquí llegas?».
\end{verse}

Ulises cuenta cómo llegó a Ogigia, la morada de Calipso, después de que Zeus hundiera su barco, cómo pasó allí siete años y cómo el viaje en balsa terminó en naufragio por intervención de Poseidón. Luego dice la verdad en lo que respecta a Nausícaa:

\begin{verse}
Se hundía el sol cuando, al fin, me soltó el dulce sueño. Mis ojos\\
vieron entonces allí a unas esclavas jugando en el soto,\\
las de tu hija, y tu hija con ellas, diosa en su coro.\\
Le supliqué; y errada no fue en su talante piadoso,\\
ese que nunca esperaras que un joven que encuentras de pronto\\
te regalara: los jóvenes siempre son faltos de aplomo.\\
Y es que ella a mí me dio pan, me dio vino ardinegro ---y a colmo---,\\
y me bañó ella en el río, y me dio el atavío que porto.\\
Hondodolido yo estoy, mas te he dicho verdad, verdad solo.
\end{verse}

A lo que responde Alcínoo:

\begin{verse}
«En algo, huésped, sí erró en aquel su talante mi hija,\\
sí, que con sus camareras aquí, a nuestra casa, tendría\\
ella que haberte traído, pues que a ella el ruego le hacías».
\end{verse}

Y Ulises:

\begin{verse}
«No me la riñas por eso, señor, a esa justa doncella,\\
que a mí ella sí me invitó a seguirla con sus asistentas,\\
pero le dije que no, por temor y también por vergüenza;\\
vete a saber si, en viéndome tú, mala mueca me hicieras,\\
que muy celosa es la raza de quienes pisamos la tierra».
\end{verse}

«Que muy celosa es la raza de quienes pisamos la tierra». Hay que decir que Odiseo ---es decir, Homero--- entiende muy bien la condición humana.

Y en este punto Alcínoo nos sale por peteneras. Recuerda, lectora de agudos ojos, lector de penetrante mirada, que todo lo que sabe el rey del mendigo que tiene delante es que se trata de un pobre hombre, un náufrago. Pero, o bien porque la referencia a Calipso le ha dado una firme pista ---pocos mortales se pasan siete años compartiendo casa y lecho con una ninfa---, o bien porque es un hombre acostumbrado a juzgar más allá de las apariencias, el rey ha decidido al vuelo que el extranjero es un buen partido y así se lo hace saber:

\begin{verse}
A lo que ya contestaba Alcínoo de nuevo diciendo:\\
«No llevo yo un corazón, extranjero, tal en el pecho\\
que se enfurruñe tan sin razón: cada cosa en su asiento.\\
Así Zeus Padre y Atena y Apolo tuvieran dispuesto\\
que siendo tú tal cual eres, y nuestros pensares gemelos,\\
a esta hija mía tomaras y fueras llamado mi yerno
\end{verse}

Pero acto seguido le dice también que, si prefiere volver a su casa, hará todo lo que esté en su mano para devolverlo a su tierra.

\begin{verse}
[...] casa y bienes pondríalos yo si, queriendo,\\
te nos quedaras; que aquí, sin querer, no te va a poner freno\\
ninguno de los feacios, que no iba a gustar eso a Zeus.\\
Tiempo sí pongo a tu guarda y escolta, y querrás bien saberlo,\\
será mañana: rendido tú entonces allí por el sueño,\\
te dormirás, y a tu patria y tu casa remando en mar quieto\\
te llevarán, o allá donde digas que tienes apego,\\
incluso si eso supone llegar hasta Eubea y más lejos,\\
la tierra más retirada nos dicen los que un día la vieron\\
de nuestros hombres, la vez en que a aquel Radamantis el güero\\
diéronle escolta para visitarse con Ticio Terrero.\\
Sí, que allí fueron y sin fatigarse en un día cumplieron,\\
y regresaron en ese día mismo a casa de nuevo.\\
Por ti sabrás tú también cuán magníficos son mis veleros\\
y esos muchachos que el vómito arrancan al mar con su remo.
\end{verse}

¡Magia de nuevo! Hasta aquí, podíamos conformarnos con la idea de que los feacios son, simplemente, buenos marinos. Pero Alcínoo le asegura a Odiseo que sus barcos pueden llevarle a Eubea ---como quien dice el fin del mundo--- en un solo día \emph{y regresar en esa misma jornada}.

Hay una correspondencia encantadora con la visita de Telémaco a Esparta. Terminada la conversación, Arete en persona se ocupa de acomodar a Odiseo ---por supuesto bajo el soportal, el mismo lugar donde Helena ha acomodado a Telémaco--- y, una vez que ha cumplido con su deber, Alcínoo y Arete hacen exactamente lo mismo que Helena y Menelao:

\begin{verse}
Alcínoo, en cambio, en lo hondo de la alta mansión se acostaba;\\
con él su esposa de ley, aliñándole el lecho, se encama.
\end{verse}


\section*{Bardos}

Si el bueno de Christopher se hubiera tomado la molestia de incluir Esqueria en su \emph{Odisea}, Travis Scott podría haberse lucido por segunda vez interpretando el rol del bardo Demódoco, en cuyos labios Homero pone la historia del caballo de Troya. La escena ocurre al día siguiente de la llegada de Ulises, durante el banquete que da Alcínoo en su honor, y es el propio héroe quien le pide que cante:

\begin{verse}
«Te doy, Demódoco, a ti sobre todos los hombres mi elogio.\\
La Musa, la hija de Zeus, debió de enseñarte, o Apolo,\\
cuando el destino de Acaya te oímos cantar de aquel modo,\\
sergas que aqueos hicieron, pasaron, sufrieron, y todo\\
como si hubieses estado tú allí o escuchado eso de otro.\\
Pero ahora cambia de tema y canta, en idéntico modo,\\
de aquel caballo que Epeo y Atena formaron con troncos\\
y que el divino Odiseo subió a la almodóvar en dolo\\
cargado con los guerreros que hicieron de Ilión un despojo.\\
Y si esto tú me lo cuentas en ley y razón, testimonio\\
ya seré yo ante los hombres por siempre de que ha sido un próvido\\
dios el que ha dado a tu canto su verbo luzdivinoso».
\end{verse}

Observa, asombrada lectora, patidifuso lector, que incluso para los estándares de Tobal, «luzdivinoso» es todo un hallazgo. Poco cuesta imaginar al bardo, encarnado por Scott, acompañándose de su lira, cantando la treta más sucia y célebre de la historia de Occidente:

\begin{verse}
Dijo; él a impulso del dios comenzó, y alumbraba su canto:\\
allá en las naves bellibancadas por fin embarcados,\\
luego de haber dado fuego a sus tiendas, se ve navegando\\
a los argivos, mas tiene el glorioso Odiseo sentados\\
otros con él ya en la plaza, ocultos en el caballo;\\
fueron los propios troyanos los que a la alcolea lo izaron.\\
Y allí él estaba, el caballo, y, sin orden ni fin perorando,\\
a su redor hacen junta los teucros; tres son sus mandados:\\
o hender el hueco madero a bronce de un fiero tajo,\\
o izarlo a cimadevilla y después desde allí despeñarlo,\\
o de él hacer un gran voto a los dioses para encandilarlos;\\
fue este mandado postrero en el que al final acabaron;\\
que era la suerte de la ciudad caer si resguardo\\
daba a ese enorme caballo de leños, cajón de esos bravos\\
argivos que a los troyanos matanza y parca llevaron.\\
Cantaba cómo los hijos de Acaya la plaza arrasaron\\
abandonando el rehueco armadijo al saltar del caballo.\\
Cómo de la alta ciudad cada cual por un lado entró a saco,\\
mientras marchose a buscar de Deífobo, aparte, el palacio,\\
a Ares gemelo, Odiseo con el pardediós Menelao.\\
Que allí la lucha más fiera, se dice, él había arrostrado\\
y la venció por la mano de Atena colosalánimo.
\end{verse}

No deja de ser curioso que Homero dedique menos versos a la escena del caballo y la posterior invasión de Troya que a comentar las riquezas del palacio de Alcínoo. Odiseo, al oírlo, no puede contener el llanto ---ya hemos dicho que los aqueos son muy llorones---:

\begin{verse}
Eso cantaba el glorioso cantor; y se desleía\\
párpado abajo Odiseo en lágrimas por sus mejillas.\\
Y como llora la esposa que a su hombre se abraza tendida\\
cuando delante él cayó de su gente mientras combatía\\
por alejar de su pueblo y sus hijos la hora negrísima,\\
y al ver que allí se le muere y que todavía palpita\\
gime en un grito sobre él derramada, y ya, con las picas\\
otros atrás golpeando sus hombros y sus costillas,\\
la llevan a esclavitud, a darle fatiga y desdichas,\\
y en el dolor lastimero marchítansele las mejillas,\\
tan lastimoso era el llanto que aquel de sus ojos vertía.\\
Mas esconder pudo a todos ese ir de sus lágrimas vivas,\\
tan solo Alcínoo lo vio y lo notó, pues tenía su silla\\
casi pegada a la de él y sus hondos suspiros sentía.
\end{verse}

Alcínoo se da cuenta de que su huésped lo está pasando mal y le pide al bardo que cese el canto:

\begin{verse}
Oíd, los que de feacios la guía tenéis y gobierno,\\
y que Demódoco enfrene aquí su sonero instrumento,\\
que es que no a todos complace el cantar que nos viene diciendo.\\
Desde el momento en que entramos en cena y se alzó el claro aedo,\\
de respirar no ha dejado en sollozos, y bien lastimeros,\\
el extranjero: sin duda ha cercado su alma algún duelo.
\end{verse}

A continuación se dirige a Ulises y le pide que cuente su nombre y su procedencia. Repara, amable lectora, gentil lector, en que a estas alturas, cuando los feacios llevan un día completo agasajando al forastero, todavía no saben cómo se llama.

\begin{verse}
Por eso tú hoy no me ocultes en trazas de buen marrullero\\
lo que de ti quiero oír: que hables tú, también es lo más bueno.\\
El nombre di que tus padres te daban allá lo primero\\
y tus vecinos y los que vecinos son de ellos linderos.\\
Que hombre que viva del todo sin nombre no existe, de cierto,\\
por vil ni noble que sea, una vez que es nacido a su tiempo;\\
a todos, sí, nos lo imponen los padres en cuanto nacemos.\\
Dime tu tierra y tu gente y también tu ciudad, porque luego\\
allí te lleven las naves guiadas por su pensamiento;
\end{verse}

Pero no es solo la curiosidad lo que mueve a Alcínoo. Necesita saber la ciudad a la que sus naves, mágicamente, van a transportar al extranjero:

\begin{verse}
no vas a ver entre los feacios ningún timonero,\\
y es que timones no hay, como sí hay en las naves del resto,\\
pues por sí solas conocen el plan que los hombres tenemos,\\
y las ciudades conocen del mundo y sus fértiles suelos,\\
y van allá de la gran hondonal de la mar casi al vuelo\\
en niebla envueltas y nubes; y en ellas no hay nunca recelo\\
de que acontezca algo malo, ni de que se pierdan hay miedo.
\end{verse}

Naves sin timonel, capaces de llegar a cualquier parte en un abrir y cerrar de ojos, siempre seguras y siempre certeras. Demasiado bonito para durar en un mundo lleno de dioses celosos y abusones. Y en efecto, los siguientes versos presagian ya una futura catástrofe:

\begin{verse}
Mas una cosa sí hay que a mi padre escuché en otro tiempo:\\
que iba a rabiar Poseidón ---siempre estaba Nausítoo diciéndolo---\\
porque les damos a todos escolta y seguro regreso.\\
Y que vendría ese día en que a un bien construido velero\\
de los feacios, volviendo de escolta en el mar neblinegro,\\
lo troncharía, y con altas montañas pondríanos cerco.
\end{verse}

Y eso es exactamente lo que va a ocurrir: la recompensa que los caprichosos dioses van a darles a los nobles feacios por cumplir a rajatabla la \emph{xenía}. Hundir sus naves y aislar para siempre Esqueria del mundo. ¿Por qué? Por una rabieta de Poseidón, que no puede soportar que, gracias a ellos, Ulises haya regresado a Ítaca.

Por el momento, sin embargo, la profecía es solo profecía y Alcínoo está más preocupado por saber a qué atenerse con su huésped.

\begin{verse}
Pero habla tú y dime a mí palabra vera de esto:\\
por dónde fuiste perdido y de hombres llegado a qué pueblos\\
y bien pobladas ciudades; y dime también si eran ellos,\\
sus habitantes, brutales, difíciles, no justicieros,\\
o acogedores y de corazón temeroso del cielo.\\
Dime también por qué lloras tan muy sollozando alma adentro\\
si oyes de Ilión y los dánaos argivos el sino funesto.\\
Eso lo obraron los dioses; la muerte de tantos la urdieron\\
para que luego sirviera de canto a los venideros.\\
¿Es que delante de Troya te vino a morir algún deudo\\
bravo, tu yerno o tu suegro, que por matrimonio son estos\\
los más cercanos de veras, de sangre después y abolengo?\\
¿O un camarada tal vez siempre en gracia contigo y apego,\\
bravo también? Inferior a un hermano de madre ---es bien cierto---\\
en nada lo es el amigo de aliento largo y discreto.
\end{verse}

Y Ulises, por fin, dice la verdad:

\begin{verse}
Empezaré por mi nombre, vosotros tenéis que saberlo\\
por si nos llega, si aquí mi hora negra yo esquivo, el momento\\
de que en mi casa yo sea, por lejos que esté, anfitrión vuestro.\\
Soy Odiseo Laertiada; y por tantas tretas que llevo\\
en boca voy de los hombres, y llega mi fama hasta el cielo.
\end{verse}

Y a continuación, empieza a contar su historia.

Permíteme, tolerante lectora, generoso lector, que abandone ahora la corte de los feacios y vuelva a Ítaca. En mi \emph{Abandonando Ítaca}\footnote{\url{https://eolasediciones.es/libro/abandonando-itaca-leaving-ithaca/}}, el anciano Ulises recuerda:


\begin{verse}
Durante muchas noches has contado tu historia,\\
en la corte del padre de Nausícaa.\\
Era un rey, pero tu memoria se niega a recordar el nombre.\\
Te pasa cada vez más a menudo, primero le pasó a Laertes,\\
olvidas el nombre de las personas,\\
ya no puedes recordar las caras,\\
ni puedes decir lo que hiciste ayer.\\
En cambio, el pasado se muestra con colores tan vivos,\\
que a menudo temes que los muertos se adueñen de los vivos.

Sí, les contaste cómo empezó la Gran Guerra.\\
Las razones fueron intrascendentes,\\
y sin embargo tan inevitables.\\
¿Cómo no creer en el destino, recordando cómo empezó todo?\\
¿Cómo no pensar que todos los hombres son juguetes de la fatalidad?
\end{verse}


\begin{verse}
Ninguno de vosotros, reyes, se preocupó por Menelao,\\
---que era más dulce con sus caballos que con su esposa---,\\
y aún menos os gustaba su hermano, el retorcido Agamenón.\\
Pero todos estabais atados por un juramento o, al menos,\\
eso dijisteis a vuestros hijos y vuestras esposas\\
y eso declarasteis en los templos, mientras vitoreaban vuestros súbditos\\
y los sacerdotes asentían con solemnidad.

El juramento era real. En la boda de Helena,\\
se pidió a los pretendientes que establecieran un vínculo irrompible,\\
con el que triunfase, ayudándole si alguna vez surgiese la necesidad.\\
Y todos os comprometisteis, poniendo por testigos a los dioses.\\
¿Cómo podías no hacerlo? Todos estabais convencidos\\
de ser el elegido. Y el trato era dulce para el ganador, desposar\\
la belleza y ganar el apoyo de los otros príncipes,\\
¿quién podría resistirse a un trato así? No tú, orgulloso Ulises,\\
el más inteligente de todos vosotros.
\end{verse}

Viejo y amargado, Ulises no ha perdido la lucidez ni la capacidad de reírse de sí mismo.

\begin{verse}
Últimamente has llegado a preguntarte\\
si eras, después de todo, tan inteligente.\\
Curioso, ahora que tus pensamientos\\
son a menudo confusos y necesitas un esfuerzo para conversar,\\
ahora que te adormeces tanto y que tu espíritu deambula con frecuencia,\\
ahora que ves en los ojos de la familia y de los sirvientes,\\
la compasión ---¿o es la condescendencia?--- hacia el anciano,\\
apenas disfrazada de respetuosa preocupación,\\
has llegado a dudar de muchas cosas\\
que eran tan obvias en tu juventud.

Y una de ellas es tu supuesta astucia.\\
Claro que eras artero, y avispado,\\
mucho más que los otros.\\
Pero ¿era un logro tan grande?\\
Al fin y al cabo, la mayoría de ellos eran montañas de carne,\\
como el bueno de Áyax, o incluso Aquiles el divino,\\
todo músculos y sin ningún ingenio. Pero, si eras tan inteligente,\\
¿cómo no entendiste que la más hermosa mujer\\
que jamás pisó la Tierra no iba a ser otorgada\\
al pobre jefe de un islote, que solo podía ofrecer\\
cabras, bravuconadas y sus trucos de sabelotodo?

Finalmente, Helena fue para el mejor\\
postor, como ha ocurrido siempre.\\
Quizás se debatiese si era una mejor opción\\
Aquiles, pero la verdad es que un gran héroe\\
no es tan buena inversión como un rey poderoso.\\
Así, Helena desposó al espartano,\\
y todos regresaron a casa sintiéndose engañados,\\
---incluso tú, que te habías forjado fama de embustero---\\
y sin ninguna intención de cumplir vuestras promesas.
\end{verse}

La competición por la mano de la bella Helena y el pacto posterior de los reyes no aparecen en Homero, sino bastante más tarde, en el \emph{Catálogo de las mujeres} atribuido a Hesíodo (siglo VI a.C.), pero, dado cómo se las gastaban los caudillos aqueos, podía haber salido de la pluma ---quiero decir el punzón--- del Poeta. La vuelta de tuerca que nos ofrece el anciano Odiseo es esta:

\begin{verse}
Solo que, cuando Agamenón llamó,\\
no solamente recordó el juramento\\
e invocó vuestro honor. No, también describió\\
una maravillosa Troya, la ciudad de altos muros\\
e incontables riquezas. Así que ayudar al amigo traicionado\\
vino con un buen extra. Después de todo, el gran ejército aqueo\\
no debería tener problemas para asaltar el lugar en pocas semanas,\\
y los líderes se repartirían el inimaginable botín.

Todos se engancharon con facilidad,\\
al ser la codicia más poderosa que la envidia y el rencor,\\
y, además, a todos os alegraba que Helena\\
hubiese hecho cornudo al soberbio espartano;\\
también os hacía felices vengar la ofensa a sangre y fuego,\\
para que vuestras esposas, vuestras hijas y hermanas,\\
tomasen buena nota de quién estaba al mando.
\end{verse}

Este Ulises tiene muy poco que ver con el santo varón que nos pinta Nolan, obligado a seguir a Agamenón para evitar represalias. ¿Cuál de los dos te parece más convincente, juiciosa lectora, sosegado lector? El bueno de Matt Damon no da abasto poniéndole excusas a Penélope:

---Tengo que ir o tomará represalias sobre mi familia.

---Pues nos vamos todos ---contesta Penélope---. En tu nave más veloz, rumbo al sol poniente.

(sigue una bella escena con la familia sobre la cubierta del raudo navío, navegando hacia el horizonte).

---No, que entonces tomaría represalias sobre mi pueblo.

¡Cuánta nobleza, que no acaba de cuadrar con el tipo que urde la treta del caballo! ¿Pero y si lo de Troya fue un mal momento, una excepción, y si, como quiere Christopher, Odiseo es un buen tipo que tuvo un mal día? Regresemos a la corte de Alcínoo y veamos lo que hace Ulises nada más zarpar de Troya.

\begin{verse}
El viento nos empujó de Ilión al lugar de los cícones,\\
a Ismaro. Allí saqueé la ciudad, muerte di a sus civiles,\\
y tras sacar de los muros mujeres, riqueza y vivires,\\
partes hicimos, que a nadie en demanda por mengua me quise.
\end{verse}

Nope, ni los diez años en las playas de Ilión ni el saqueo y destrucción de Troya han servido para que Ulises deje de ser un forajido. Si el incendio de la ciudad sagrada le ha producido algún trauma, eso no le impide entrar a sangre y fuego en el primer poblado que encuentra en su camino de vuelta. Podría argumentarse que lo hace porque los cícones eran aliados de los troyanos, pero la guerra ya ha terminado y, en todo caso, Ulises no nos ofrece razón alguna que justifique su conducta. Declara que entró a saquear, matar, incendiar y robar mujeres, como quien dice que fue a buscar tabaco.

De vuelta en su isla, al anciano no le queda otra que reconocer la verdad:


\begin{verse}
Al menos ahora eres lo bastante lúcido\\
para no engañarte, como solías hacer\\
cuando te imaginabas ser el gran Ulises,\\
el héroe y el viajero y el amante de Circe.\\
Cuando recuerdas, todo lo que ves\\
es ruido y furia. Un matón cruel y astuto\\
que tuvo más suerte que los demás.
\end{verse}


\begin{verse}
Tus manos aún estaban mojadas con sangre de troyanos,\\
cuando tú y tus hombres atacabais a los cícones\\
en vuestro camino de regreso a casa. No tenías nada contra ellos,\\
y tus barcos ya estaban llenos del botín\\
de Ilión. Y sin embargo, no podías perder la oportunidad\\
de saquear otra ciudad. ¿Con qué propósito?\\
Los cícones se defendieron con bravura, obligándoos\\
a volver al mar. Sí, robaste más joyas y más oro,\\
pero el oro no es comestible y las joyas no sacian la sed.\\
Peor aún, tus barcos estaban tan sobrecargados\\
que te arriesgaste a naufragar. ¡Solo por avaricia! Pero en ese momento\\
no viste la ironía; estabas demasiado furioso,\\
lanzando a las olas tu tesoro.

Quienes cantan tus hazañas cuentan que unas ninfas vinieron,\\
pidiendo más, e hicieron enfadar a Poseidón.\\
Ninfas que tú no viste, y a medida que tu fe en los dioses disminuye,\\
has llegado a comprender que los hombres inventaron a los inmortales,\\
para dar un sentido a sus propios errores.
\end{verse}























\end{document}
